\documentclass[12pt,a4paper,titlepage]{article}
%va sempre messo article per "program documentation"
\usepackage[italian]{babel}
\usepackage[T1]{fontenc}
\usepackage[latin1]{inputenc}
\usepackage{titlesec}
\usepackage{hyperref}
\usepackage[a4paper,top=2cm,bottom=2cm,left=1cm,right=1cm]{geometry}
\usepackage{soulutf8,color}
\usepackage{multirow}
\usepackage{eurosym} %per il simbolo dell' \euro
\usepackage{fancyhdr} %intestazioni
\usepackage{graphicx}


\usepackage[acronym]{glossaries}
\newglossaryentry{parola}{name=parola,description={Si spiega da sè}}
\newacronym{guit}{GuIT{}}{Gruppo Utilizzatori Italiani di \LaTeX}
\addto\captionsitalian{\renewcommand{\glossaryname}{Glossario}}
\makeglossaries

%----------------------------------------------------------- INIZIO TEMPLATE TITOLO

\usepackage{xcolor} % Importa i colori per la prima pagina
\usepackage{fix-cm} % Permette l'incremento del font oltre misura


\newcommand{\HRule}[1]{\hfill \rule{0.2\linewidth}{#1}} % Horizontal rule at the bottom of the page, adjust width here

\definecolor{grey}{rgb}{0.9,0.9,0.9} % Colore del box del titolo

\begin{document}
	
	\thispagestyle{empty} % Toglie il numero della pagina nella prima pagina
	
	%----------------------------------------------------------------------------------------
	%	TITLE SECTION
	%----------------------------------------------------------------------------------------
	
	\colorbox{grey}{
		\parbox[t]{1.0\linewidth}{
			\centering \fontsize{50pt}{80pt}\selectfont % Il primo è la grandezza del font, il secondo lo spazio lasciato
			\vspace*{0.7cm} % Spazio dall'inizio del box al titolo
		
				\raggedleft
				\includegraphics[width=0.7\linewidth]{../../LogoSWEgGroupSFONDOVUOTO}
			
			\hfill Glossario \\
			
			\vspace*{0.7cm} % Spazio dalla fine del testo alla fine del box
		}
	}
	
	%----------------------------------------------------------------------------------------
	
	\vfill % Spazio dalla fine del box alle altre informazioni
	
	%----------------------------------------------------------------------------------------
	%	Informazioni sul documento
	%----------------------------------------------------------------------------------------
	
	{\centering \large 
		\hfill \textbf{Versione} 1.0.0 \\
		\hfill \textbf{Redazione} Sebastiano Marchesini \\
		\hfill \textbf{Verifica} Piergiorgio Danieli \\
		\hfill \textbf{Responsabile} Sebastiano Marchesini \\
		\hfill \textbf{Uso} Interno \\
		\hfill \textbf{Destinato} SWEg Group \\ 
		
		\HRule{1pt}
	
	\textbf{Sommario} \\
	Testo destinato alla comprensione di sigle, inglesismi e termini tecnici all'interno dei documenti.

	} % Linea orizzontale di estetica


	%----------------------------------------------------------------------------------------
	
	\clearpage % Parta bianca finale della pagina

%----------------------------------------------------------- FINE TEMPLATE TITOLO

\lhead{\includegraphics[width=0.2\linewidth]{../../LogoSWEgGroup}}
\chead{}
\lfoot{Glossario}
\cfoot{}
\rfoot{\thepage}
\renewcommand{\headrulewidth}{0.2pt}
\renewcommand{\footrulewidth}{0.2pt}

\tableofcontents
%crea indice automaticamente
\thispagestyle{empty}

\newpage


\rhead{Introduzione}
\section{Introduzione}
\subsection{Scopo del Documento}
Un glossario è una raccolta di termini di un ambito specifico e circoscritto. In alcuni casi il termine glossario è utilizzato per indicare la raccolta e relativa spiegazione dei termini desueti o specialistici utilizzati in un libro (generalmente collocato alla fine del testo), nel nostro caso la documentazione. \\
 
\subsection{Riferimenti}
Come riferimenti per i termini del glossario si è utilizzato
\begin{itemize}
	\item \textbf{Wikipedia}: \\
	\textcolor{blue}{\url{https://it.wikipedia.org/}}.
	\item \textbf{Dizionario Treccani}: \\
	\textcolor{blue}{\url{http://www.treccani.it/vocabolario/}}.
	\item \textbf{Dispense fornite dal corso Ingegneria del Software 2016}: \\
	\textcolor{blue}{\url{http://www.math.unipd.it/~tullio/IS-1/2016/Dispense}}.
		\item \textbf{Siti ufficiali delle software-house}: \\
\end{itemize}

\newpage	

\rhead{A}	
\section{A}
\begin{trivlist}
\item \textbf{Account} : Il complesso dei dati identificativi di un utente, che gli consentono l'accesso a un servizio telematico.
\item \textbf{Amministratore} : È responsabile dell'efficienza e dell'operatività dell'ambiente di sviluppo, della redazione e attuazione di piani e procedure di Gestione per la Qualità. Controlla versioni e configurazioni del prodotto. Gestisce l'archivio della documentazione di progetto. Collabora alla redazione del Piano di Progetto. Redige le Norme di Progetto per conto del Responsabile.
\item \textbf{Analisi Dinamica} : Processo di valutazione di un sistema software o di un suo componente basato sulla osservazione del suo comportamento in esecuzione.
\item \textbf{Analisi Statica} : Non richiede l'esecuzione del prodotto. Studia le caratteristiche del codice sorgente e della documentazione
\item \textbf{Analista} : È responsabile delle attività di analisi. Redige lo Studio di Fattibilità (documento interno al gruppo) e l'Analisi dei Requisiti.
\item \textbf{API} : È l'acronimo di Application Programming Interface, nonché di Application Program Interface.
\item \textbf{APIkey} : È una chiave generata dal sistema che permette ad un utente di utilizzare una determinata API.
\item \textbf{APIMarket (o APIM)} :È l'acronimo di API Market ed è anche l'acronimo titolo del capitolato scelto.
\item \textbf{Asana} : Asana è una applicazione web e mobile progettata per aiutare i team monitorare il loro lavoro.
\item \textbf{Astah} : Astah è uno strumento di progettazione del software che supporta UML, ERD, DFD, Diagramma di Flusso, Mindmap, SysML, GSN e altri.
\item \textbf{Attori} : Rappresenta un ruolo coperto da un certo insieme di entità interagenti col sistema (inclusi utenti umani, altri sistemi software, dispositivi hardware e così via). Un ruolo corrisponde a una certa famiglia di interazioni correlate che l'attore intraprende col sistema. Usato nei diagrammi UML.
\end{trivlist}

\rhead{B}	
\section{B}
\begin{trivlist}
\item \textbf{Back End (BE)} : Parte del sistema che permette l'effettivo funzionamento delle iterazioni con l'utente.
\item \textbf{Backlog} : Operazione in attesa di essere eseguita. Prodotta dal Project Manager.
\item \textbf{Bisogni Espliciti} : Derivanti dal capitolato d'appalto sono effitavamente i bisogni del cliente.
\item \textbf{Bisogni Impliciti} : Derivati dal dominio sono bisogni tratti da standard di qualità e altro.
\item \textbf{Branch} : Comando di Git adibito a creare un nuovo ramo. Dove è possibile fare modifiche senza modificare il ramo principale: master.
\item \textbf{Browser} : Particolare programma per navigare in Internet che inoltra la richiesta di un documento alla rete e ne consente la visualizzazione una volta arrivato.
\item \textbf{Budget} : Disponibilità finanziaria. Bilancio preventivato.
\end{trivlist}

\rhead{C}	
\section{C}
\begin{trivlist}
\item \textbf{Capitolato D'Appalto} : Il capitolato è un documento tecnico che vi fa riferimento per definire in quella sede le specifiche tecniche delle opere che andranno ad eseguirsi per effetto del contratto stesso, di cui è solitamente parte integrante.
\item \textbf{Cascata} : Il modello a cascata (waterfall model in inglese) o ciclo di vita a cascata (waterfall lifecycle) è il più tradizionale modello di ciclo di vita del software. Le fasi sono sequenziali e distinte.
\item \textbf{Collaudo} : Controllo, validazione dei requisiti tecnologici del progetto in rapporto a una tabella di caratteristiche singolarmente o universalmente prestabilite.
\item \textbf{Committente} : È la figura che commissiona un lavoro, indipendentemente dall'entità o dall'importo.
\item \textbf{Crittografia} : È la branca della crittologia che tratta delle "scritture nascoste", ovvero dei metodi per rendere un messaggio "offuscato" in modo da non essere comprensibile/intelligibile a persone non autorizzate a leggerlo.
\item \textbf{CSS} : Un linguaggio usato per definire la formattazione di documenti HTML, XHTML e XML ad esempio i siti web e relative pagine web. Di solito usato per gestire l'interfaccia di visualizzazione dell'utente e dispositivi.
\end{trivlist}

\rhead{D}	
\section{D}
\begin{trivlist}
\item \textbf{Dati Tecnici} : In riferimento alle API i dati tecnici sono i dati prelevati dall'interfaccia gateway per poi essere catalogati e confrontati.
\item \textbf{Diagramma dei Casi D'Uso (UML)} : Sono diagrammi dedicati alla descrizione delle funzioni o servizi offerti da un sistema, così come sono percepiti e utilizzati dagli attori che interagiscono col sistema stesso. Sono impiegati soprattutto nel contesto della Use Case View (vista dei casi d'uso) di un modello, e in tal caso si possono considerare come uno strumento di rappresentazione dei requisiti funzionali di un sistema.
\item \textbf{Documentazione} : L'insieme dei documenti da presentare insieme al progetto fisico. Comprendono in particolare : Analisi dei Requisiti, Piano di Qualifica, Piano di Progetto, Norme di Progetto.
\end{trivlist}

\rhead{E}	
\section{E}
\begin{trivlist}
\item \textbf{Efficaci} : Conformi alle attese.
\item \textbf{Efficienti} : Contenimento dei consumi per raggiungere un obiettivo (diminuire costi e tempi).
\end{trivlist}

\rhead{F}	
\section{F}
\begin{trivlist}
\item \textbf{File System} : Indica informalmente un meccanismo con il quale i file sono posizionati ed organizzati o su un dispositivo di archiviazione o su una memoria di massa, come un disco rigido o un CD-ROM e, in casi eccezionali, anche sulla RAM.
\item \textbf{Fork} : Indica lo sviluppo di un nuovo progetto software che parte dal codice sorgente di un altro già esistente, ad opera di un programmatore.
\item \textbf{Fornitore} : Nel nostro caso il fornitore è il professore che si pone da tramite tra noi e l'azienda.
\item \textbf{Framework} : Piattaforma che funge da strato intermedio tra un sistema operativo e il software che lo utilizza.
\item \textbf{Front End (FE)} : La parte visibile all'utente e con cui egli può interagire (interfaccia utente).
\end{trivlist}

\rhead{G}	
\section{G}
\begin{trivlist}
\item \textbf{Git} : Git è un software di controllo versione distribuito utilizzabile da interfaccia a riga di comando, creato da Linus Torvalds nel 2005.
\item \textbf{Google Play} : Google Play, detto anche Google Play Store, è un negozio virtuale online di applicazioni, brani musicali, pellicole cinematografiche, libri e riviste sviluppato da Google Inc. principalmente per offrire servizi ai dispositivi mobili Android.
\item \textbf{GUI} : Acronimo di Graphic User Interface, interfaccia grafica costituita da icone, finestre e tasti selezionabili attraverso il mouse.
\item \textbf{Gulpease} : L'indice Gulpease è un indice di leggibilità di un testo tarato sulla lingua italiana. Rispetto ad altri ha il vantaggio di utilizzare la lunghezza delle parole in lettere anziché in sillabe, semplificandone il calcolo automatico.
\end{trivlist}


\rhead{H}	
\section{H}
\begin{trivlist}
	\item \textbf{HTML} : Acronimo di Hyper Text Markup Language, è un linguaggio di markup utilizzato per la creazione di documenti destinati al web. La visualizzazione dei file HTML, chiamati pagine web, è affidata ad un browser che traduce il linguaggio HTML in contenuti consultabili dagli utenti.
\end{trivlist}

\rhead{I}	
\section{I}
\begin{trivlist}
\item \textbf{IEC} : Acronimo di  International Electrotechnical Commission,  è un'organizzazione internazionale per la definizione di standard in materia di elettricità, elettronica e tecnologie correlate.
\item \textbf{Inspection} : Parte dell'analisi statica modello che cerca gli errori più probabili tramite check-list . C'è bisogno della definizione di una lista di controllo.
\item \textbf{Instagantt} : È una web app per la creazione di diagrammi di Gantt che aiutano visivamente intendono pianificazioni dei progetti con il team.
\item \textbf{ISO} : Primo elemento di composti, derivati dal greco o formati modernamente, nei quali indica ‘uguaglianza’ o anche ‘affinità’. Nel nostro caso abbinato al termine IEC indica uno standard.
\item \textbf{ItalianaSoftware} : ItalianaSoftware è la società che ha ideato Jolie. Fornisce soluzioni microservizi e prodotti come strumenti amministrazione microservizi e SAP
\end{trivlist}

\rhead{J}	
\section{J}
\begin{trivlist}
\item \textbf{Javascript} : È un linguaggio di scripting orientato agli oggetti e agli eventi, comunemente utilizzato nella programmazione Web lato client per la creazione, in siti web e applicazioni web, di effetti dinamici interattivi tramite funzioni di script invocate da eventi innescati a loro volta in vari modi dall'utente sulla pagina web in uso (mouse, tastiera, caricamento della pagina ecc...).
\item \textbf{Jolie} : È un linguaggio open source simile a Java e C. Con lo scopo di facilitare la progettazione dei microservizi.
\end{trivlist}

\rhead{L}	
\section{L}
\begin{trivlist}
\item \textbf{Latex} : È un linguaggio di markup usato per la preparazione di testi basato sul programma di composizione tipografica TEX.
\item \textbf{Libreria} :  Si tratta di un insieme di funzioni o strutture dati predefinite e predisposte per essere collegate ad un programma software attraverso opportuno collegamento. 
\item \textbf{Linux} : È una famiglia di sistemi operativi di tipo Unix-like, rilasciati sotto varie possibili distribuzioni, aventi la caratteristica comune di utilizzare come nucleo Linux.
\end{trivlist}

\rhead{M}	
\section{M}
\begin{trivlist}
\item \textbf{Mac Os X} : È il sistema operativo sviluppato da Apple Inc. per i computer Macintosh.
\item \textbf{Merge} : Indica una operazione di costruzione di un unico file dati, a partire da file separati.
\item \textbf{Microservizio} : Nati essenzialmente come soluzione al problema dell'integrazione tra macrosistemi operanti su piattaforme differenti; sono nati cioè con l'obiettivo di fare inter-operare in modo organico e non occasionale sistemi tra loro differenti. 
\item \textbf{Milestone} : Indica importanti traguardi intermedi nello svolgimento del progetto.
\item \textbf{Misurabile} : Che può essere misurato, suscettibile di valutazione quantitativa in rapporto ad una unità di misura.
\item \textbf{Modello Agile} : Metodo di sviluppo software che propone un approccio meno strutturato e focalizzato sull'obiettivo di consegnare al cliente, in tempi brevi e frequentemente, software funzionante e di qualità.
\item \textbf{Monitoramento} : Rilevazione o verifica effettuata con l'uso di uno strumento specifico.
\item \textbf{Monolite} : Vecchio sistema e/o programma realizzato con molte dipendenze e poco flessibile.
\item \textbf{MySQL} : È il più diffuso database Open Source basato sul linguaggio SQL. MySQL è un RDBMS, ossia un sistema di gestione per database relazionali.
\end{trivlist}

\rhead{P}	
\section{P}
\begin{trivlist}
\item \textbf{Package} : Pacchetto. Di solito usato come termine specifico nel linguaggio Java.
\item \textbf{Password} : Parola o sigla di riconoscimento fornita dall'utente all'elaboratore per poter accedere a un sistema operativo, a un programma o a un file.
\item \textbf{PDCA} : Ciclo di PDCA  (plan–do–check–act) è un metodo di gestione in quattro fasi iterativo, utilizzato in attività per il controllo e il miglioramento continuo dei processi e dei prodotti.
\item \textbf{Piattaforma} : Un ambiente di esecuzione che comprende hardware e sistema operativo ed eventualmente elementi middleware specifici, application server ed altri strumenti di supporto all'esecuzione di programmi.
\item \textbf{Policy} : Nel nostro caso si intende la politica secondo la quale un utente può acquistare una API (ad esempio per un periodo di tempo, per una certa quantità di traffico dati ecc.).
\item \textbf{Progettista} : È responsabile delle attività di progettazione. Redige Specifica Tecnica, Definizione di Prodotto e la parte programmatica del Piano di Qualifica.
\item \textbf{Programmatore} : È responsabile delle attività di codifica miranti alla realizzazione del prodotto e delle componenti di ausilio necessarie per l'esecuzione delle prove di verifica e validazione.
\item \textbf{Project Manager} : O Responsabile di Progetto. È il responsabile ultimo, per conto del suo gruppo, dei risultati del progetto. Elabora ed emana piani e scadenze. Approva l'emissione di documenti. Coordina le attività del gruppo, si relaziona con il controllo di qualità interno al progetto. Redige Organigramma e Piano di Progetto. Approva l'Offerta e i relativi allegati.
\item \textbf{Proponente} : In genere l'azienda che propone il capitolato. Nel nostro caso Italiana Software.
\end{trivlist}

\rhead{R}	
\section{R}
\begin{trivlist}
\item \textbf{Raw} : Crudo. Si usa nel contesto di documento o file crudo , con errori grossolani quindi.
\item \textbf{Repository} : È un ambiente di un sistema informativo, in cui vengono gestiti i metadati. Nel nostro contesto si utilizza spesso per indicare lo spazio GIT.
\item \textbf{REST} : REpresentational State Transfer (REST) è un tipo di architettura software per i sistemi di ipertesto distribuiti come il World Wide Web.
\item \textbf{Revisione} : Esame o controllo, per lo più periodico, inteso a verificare il grado dell'efficienza, della funzionalità, della corrispondenza a determinati requisiti, in quanto può implicare apporto di modifiche o di correzioni.
\item \textbf{Revisione dei Requisiti (RR)} : Primo processo manuale che coinvolge team misti. Comprende tra i documenti singolari Studio di Fattibilità e Analisi dei Requisiti.
\item \textbf{Revisione di Accettazione (RA)} : Collaudo del sistema per accettazione della parte committente. Accertamento del soddisfacimento di tutti i requisiti utente pattutiti in RR.
\item \textbf{Revisione di Progettazione (RP)} : Può essere di due tipi : Minima (Min) e Massima (Max). Rispettivamente accertano la realizzabilità e si accordano sulle caratteristiche del prodotto da realizzare. I documenti da presentare specifici sono: nella minima la Specifica Tecnica, mentre nella massima la Definizione del Prodotto
\item \textbf{Revisione di Qualifica (RQ)} : Revisione dell'esito finale di qualifiche delle verifiche e attivazione di validazione.
\item \textbf{Ritorno Incrementale} : È un modello di sviluppo di un progetto software basato sulla successione di passi. Non vi sono iterazioni e lo si usa quando si ha una chiara visione dell'intero progetto. L'uso di questo modello di sviluppo favorisce la creazione di prototipi, ovvero parti di applicazione funzionanti, che a loro volta favoriscono il dialogo con il cliente e la validazione dei requisiti.  
\end{trivlist}

\rhead{S}	
\section{S}
\begin{trivlist}
\item \textbf{Script} : "Simile ad un programma normale ma con una complessità relativamente bassa. Si integra in un processo di configurazione automatica del sistema, in una pagina web, o comunque per svolgere mansioni accessorie e molto specifiche con
una certa linearità. Manca di una propria interfaccia grafica e richiama altri programmi per svolgere operazioni più sofisticate."
\item \textbf{Scrum} : I Team Scrum sono auto-organizzati e cross-funzionali: scelgono come meglio compiere il lavoro organizzandosi e coordinandosi al proprio interno e hanno tutte le competenze necessarie per realizzare il lavoro senza dover dipendere da nessuno al di fuori del team. Il modello di team in Scrum è progettato per ottimizzare la flessibilità, la creatività e la produttività. I Team Scrum rilasciano i prodotti in modo iterativo e incrementale, massimizzando le opportunità di feedback. I rilasci incrementali di prodotto garantiscono che una versione potenzialmente utile del prodotto funzionante sia sempre disponibile.
\item \textbf{Sistema Operativo} : È un insieme di componenti software, che rende operativi (da cui il nome) computer, apparati e dispositivi informatici.
\item \textbf{Smartphone} : Apparecchio elettronico che combina le funzioni di un telefono cellulare e di un computer palmare.
\item \textbf{Snapshot} : Istantanea. Termine utilizzato nei programmi di grafica. Rappresenta una fase di lavoro che si vuole "fermare" nel caso che le variazioni che si stanno per compiere non ci soddisfino. Le snapshot consentono quindi di vedere versioni alternative di una stessa immagine, per poter scegliere la migliore.
\item \textbf{SQA} : Acronimo di Software Quality Assurance, costiyuisce un mezzo per monitorare i processi di ingegneria del software e dei metodi utilizzati per garantire la qualità. I metodi con cui questo viene realizzato sono molteplici, e possono includere garantire la conformità ad uno o più standard, come ISO 9000 o un modello come CMMI.
\item \textbf{Statistiche} : Sono dati calcolati su un utente della piattaforma. Servono a rendere confrontabili e misurabili le caratteristice dell'utente per una visione social e più amministrabile dell'infrastruttura.
\item \textbf{Steam} : Piattaforma online molto simile ad un mercato di videogiochi dove è possibile interagire con altri utenti, acquistare giochi e caricarne propri.
\item \textbf{Sub-Task} : Sotto attività.
\end{trivlist}

\rhead{T}	
\section{T}
\begin{trivlist}
\item \textbf{Tablet} : Computer portatile di dimensioni ridotte, sul cui schermo è possibile scrivere o impartire comandi col tocco delle dita o mediante un apposito stilo.
\item \textbf{Tag} : Sequenza di caratteri con cui si marcano gli elementi di un file per successive elaborazioni.
\item \textbf{Task} : Attività.
\item \textbf{Team} : Gruppo di lavoro.
\item \textbf{Telegram} : Programma di messaggistica istantanea che si basa su numero di telefono e/o id. Può essere utilizzata su piattaforma mobile e non. 
\item \textbf{Template} : Modello predefinito che consente di creare o inserire contenuti di diverso tipo in un documento.
\item \textbf{Test} : Collaudo del software (detto anche testing o software testing secondo le denominazioni inglesi) è un procedimento, che fa parte del ciclo di vita del software, utilizzato per individuare le carenze di correttezza, completezza e affidabilità delle componenti software in corso di sviluppo.
\item \textbf{Test di Integrazione} : Test che mira alla correttezza delle interfacce.
\item \textbf{Test di Sistema} : Test di più alto livello che mira all'affidabilità, sicurezza e alle prestazioni.
\item \textbf{Test di Unità} : Test di più basso livello che mirano alla correttezza delle interfacce.
\item \textbf{TexStudio} : È una piattaforma di sviluppo per il linguaggio LaTex che agevola il programmatore a stilare la documentazione.
\item \textbf{Ticket (Ticketing)} : Un trouble ticket indica letteralmente una richiesta di assistenza, tracciata da un sistema informatico di gestione delle richieste di assistenza, che per metonimia viene indicato con lo stesso termine.
\end{trivlist}

\rhead{U}	
\section{U}
\begin{trivlist}
\item \textbf{Use Case (UC)} : Definizione inglese per idicare i Casi D'Uso. UCD non indicato nel progetto indica gli Use Case Diagram.
\item \textbf{Username} : Numero o parola utilizzati da un utente per farsi identificare da un sistema operativo, da un elaboratore o da un servizio online.
\end{trivlist}

\rhead{V}	
\section{V}
\begin{trivlist}
\item \textbf{Verificatore} : È responsabile delle attività di verifica. Redige la parte retrospettiva del Piano di Qualifica che illustra l'esito e la completezza delle verifiche e delle prove effettuate secondo il piano.
\item \textbf{Versionamento} : l controllo versione (versioning) è la gestione di versioni multiple di un insieme di informazioni. Gli strumenti software per il controllo versione sono ritenuti molto spesso necessari per la maggior parte dei progetti di sviluppo software.
\end{trivlist}

\rhead{W}	
\section{W}
\begin{trivlist}
\item \textbf{Walkthrough} : Medoto per rivelare la presenza di difetti. Più lento di Ispector, esegue una lettura critica (a largo spettro) del codice. Percorrere il codice simulandone possibili esecuzioni o il documento completamente.
\item \textbf{Windows} : Microsoft Windows (abbreviazioni comunemente utilizzate: "Windows" o "Win") è una famiglia di ambienti operativi e sistemi operativi dedicati ai personal computer, alle workstation, ai server e agli smartphone.
\item \textbf{Workspace} : Ambiente di lavoro.
\end{trivlist}

\end{document}

